\documentclass[10pt,a4paper]{amsart}
\usepackage[margin=19mm]{geometry}
\usepackage{amsmath,amssymb,mathtools}
\usepackage[scr=rsfs,cal=euler]{mathalfa}
\usepackage{xfrac}
\usepackage[unicode,hidelinks,bookmarks=false]{hyperref}
\newcommand{\ie}{i.e.\ }
\newcommand{\cf}{cf.\ }
\newcommand{\resp}{respectively\ }
\newcommand{\Subsection}{\subsection}
\providecommand{\nobreakdash}{\nobreak\hbox{-}}
\setlength{\emergencystretch}{2em}
\multlinegap0pt
\usepackage{mathtools}
\usepackage{booktabs}
\usepackage{enumitem}
\usepackage{url}
\newtheorem{theorem}{Theorem}[section]
\newtheorem{proposition}[theorem]{Proposition}
\newtheorem{lemma}[theorem]{Lemma}
\newtheorem{corollary}[theorem]{Corollary}
\newtheorem{remark}[theorem]{Remark}
\newcommand{\F}{\mathbb F}
\newcommand{\Aut}{\operatorname{Aut}}
\newcommand{\Syl}{\operatorname{Syl}}
\newcommand{\PP}{\mathbb P}
\newcommand{\cH}{\mathcal H}
\begin{document}
\title[Transitive automorphism groups of maximal curves]{Transitive automorphism groups of maximal curves}
\author{Saeed Tafazolian}
\date{}
\maketitle
\noindent\emph{Source and licence.} Transitive automorphism groups of maximal curves. Version arXiv:2609.17611v1 (CC BY 4.0). This material is distributed under the Creative Commons Attribution 4.0 International licence, http://creativecommons.org/licenses/by/4.0/, as recorded in the arXiv OAI licence metadata for this exact version and shown on the abstract page https://arxiv.org/abs/2609.17611v1. Full rights notice: licenses/arxiv-2609.17611-CC-BY-4.0.txt.\par\smallskip
\noindent\textit{Editorial scope.} Counted unit: the original \texttt{theorem} environment \texttt{thm:GMP-TI} at source lines 239--260 together with its complete proof at lines 262--301; the delivered document keeps source lines 145--369 so that every referenced label and every citation resolves inside it. Native editable source is the paper's own concatenated TeX; the AMS wrapper is editorial formatting only. No publisher class, upstream script or unrelated artwork is redistributed.
\section{Preliminaries}

We collect the facts used in the proof.

\subsection{Genus bounds}

We shall repeatedly use the upper part of the genus spectrum of a
maximal curve.  The maximal-genus characterization in part \textup{(i)}
is due to R\"uck--Stichtenoth, while the second-genus bound in part
\textup{(ii)} is due to Fuhrmann--Garcia--Torres; we use the formulations
recalled in \cite[Lemma~2.12 and Corollary~2.13]{KT}.  The equality case
in part \textup{(iii)} is \cite[Theorem~3.1]{FGT}, and the third genus
bound in part \textup{(iv)} is \cite[Corollary~3.3]{KT}.

\begin{lemma}\label{lem:gaps}
Let $X/\F_{q^2}$ be a maximal curve of genus $g$.
Then:
\begin{enumerate}[label=\textup{(\roman*)}]
\item
$g\le q(q-1)/2$, with equality if and only if $X$ is Hermitian.

\item If $X$ is not Hermitian, then
\[
g\le \frac{(q-1)^2}{4}.
\]

\item If $q$ is odd and $g=(q-1)^2/4$, then
\[
X\simeq_{\F_{q^2}}\{\,y^{(q+1)/2}=x^q+x\,\}.
\]

\item If
\[
g<\left\lfloor\frac{(q-1)^2}{4}\right\rfloor,
\]
then
\[
g\le\left\lfloor\frac{q^2-q+4}{6}\right\rfloor.
\]
\end{enumerate}
\end{lemma}

A maximal curve is supersingular.  We shall also use the following
standard consequence of the Kleiman--Serre covering result; see
\cite[Proposition~6]{Lachaud}.

\begin{lemma}\label{lem:quotientmax}
Let $X/\F_{q^2}$ be maximal, and let $H\leq\Aut(X)$ be defined over
$\F_{q^2}$.  Then the quotient $X/H$ is maximal over $\F_{q^2}$.
In particular, every positive-genus quotient of $X$ is supersingular.
\end{lemma}

\subsection{Group actions and ramification notation}

Let a finite group $G$ act on a curve $X$.  For $P\in X$, we write
\[
        G_P=\{\sigma\in G:\sigma(P)=P\}
\]
for the stabilizer of $P$, and
\[
        G\cdot P=\{\sigma(P):\sigma\in G\}
\]
for its orbit.  The orbit--stabilizer formula gives
\[
        |G|=|G_P|\,|G\cdot P|.
\]
An orbit is called short if its length is smaller than $|G|$, or
equivalently if its point stabilizers are nontrivial.

The action determines the quotient Galois cover
\[
        X\longrightarrow X/G.
\]
For $P\in X$, the ramification index at $P$ is $|G_P|$.  We write
\[
        G_P=G_P^{(0)}\supseteq G_P^{(1)}
             \supseteq G_P^{(2)}\supseteq\cdots
\]
for the lower ramification groups.  If the action is tame, then
$G_P^{(1)}=1$ and $G_P$ is cyclic.  In Section~\ref{sec:tame} we shall
call a subgroup $G\le\Aut(X)$ \emph{tame} when $p\nmid|G|$.  This is
the global hypothesis used there; in particular every inertia group of
the action is tame.

We write $\Syl_p(G)$ for the set of Sylow $p$-subgroups of $G$,
$N_G(S)$ for the normalizer of a subgroup $S\le G$, and
$O_{p'}(G)$ for the largest normal subgroup of $G$ of order prime to
$p$.  A Sylow $p$-subgroup is called TI if distinct conjugates have
trivial intersection.

We shall use the following structural result in the wild case. It is
quoted here in the geometric form needed below, in order to make clear
both its hypotheses and its $2$-transitivity conclusion.


\begin{theorem}[Guralnick--Malmskog--Pries in the form used here]\label{thm:GMP-TI}
Let $A$ act faithfully on a curve $Y$ of genus at least two over an
algebraically closed field of characteristic $p$, and let
$Q\in\Syl_p(A)$. Suppose that $Q$ is not cyclic (and, when $p=2$, not
generalized quaternion), that $Q$ fixes exactly one point of $Y$ and
acts freely on its complement, and that
\[
        N_A(Q)=QC
\]
with $C$ cyclic. If $N_A(Q)\ne A$ and $M=Z(N_A(Q))$, then $M$ is a
normal $p'$-subgroup of $A$, the quotient $A/M$ is almost simple, and
its socle is isomorphic to one of
\[
 PSL(2,p^a)\quad(a\ge2),\qquad
 PSU(3,p^a)\quad(p^a>2),
\]
\[
 Sz(2^{2a+1})\quad(p=2,\ a>1),\qquad
 {}^2G_2(3^{2a+1})'\quad(p=3).
\]
Moreover, $A$ acts doubly transitively on $\Syl_p(A)$.
\end{theorem}

\begin{proof}[Justification of the quoted form]
The normality of $M$ and the almost-simple conclusion are precisely the
geometric conclusion of \cite[Theorem~3.19]{GMP}; we record why the
$2$-transitivity assertion from \cite[Theorem~3.16]{GMP}, on which that
theorem is based, applies in the present setting.

First, $Q$ is a TI subgroup. Let $P$ be the unique fixed point of $Q$
and take $a\notin N_A(Q)$. If $Q$ and $Q^a$ fixed the same point, then
both would be Sylow $p$-subgroups of the corresponding point
stabilizer. The first ramification group is the unique Sylow
$p$-subgroup of a point stabilizer, so this would force $Q=Q^a$, a
contradiction. Thus the fixed points of $Q$ and $Q^a$ are distinct. A
nontrivial element of $Q\cap Q^a$ would fix both of them, contrary to
the assumed free action away from the unique fixed point of $Q$. Hence
$Q$ is TI.

There is no mismatch between the hypotheses used here and those in
\cite{GMP}. In Notation~3.18 of that paper the condition that $Q$ be
noncyclic (and, for $p=2$, not generalized quaternion) is explicitly
noted to be equivalent to the existence of an elementary abelian
subgroup of order $p^2$. Moreover, Theorem~3.19 assumes
$A\ne I=N_A(Q)$, which is exactly our hypothesis $N_A(Q)\ne A$.
The group-theoretic Theorem~3.16 is stated under the weaker condition
$Q\ne A$, and this is automatic here.

For completeness, we make explicit the initial reduction in the proof
of \cite[Theorem~3.16]{GMP}. Since $N_A(Q)\ne A$, choose
$a\notin N_A(Q)$. The normal $p$-subgroup $O_p(A)$ is contained in
every Sylow $p$-subgroup of $A$, and therefore
\[
        O_p(A)\le Q\cap Q^a=1.
\]
Thus $O_p(A)=1$. The argument of \cite[Theorem~3.16]{GMP} then passes
to $A/O_{p'}(A)$, applies the finite-simple-group classification quoted
there, and obtains the displayed socle list and the doubly transitive
action on the Sylow $p$-subgroups. In fact, that theorem proves the
slightly stronger identity $M=O_{p'}(A)$. The classification of finite
simple groups enters only through this step; see
\cite[Remark~3.17]{GMP}.
\end{proof}


\subsection{Tame specialization and triangle actions}\label{subsec:tame-specialization}

Suppose that $p\nmid|G|$ and that the quotient $X/G$ is rational with
exactly three branch points.  If the ramification indices are
$(a,b,c)$, with
\[
       2\le a\le b\le c,\qquad
       \frac1a+\frac1b+\frac1c<1,
\]
then Grothendieck's specialization theorem for the prime-to-$p$
fundamental group of the punctured line gives a generating triple of $G$
of orders $a,b,c$ whose product is $1$; see
\cite[Expos\'e~XIII, Corollaire~2.12]{SGA1}.  Equivalently, $G$ is a
finite quotient of the hyperbolic triangle group $\Delta(a,b,c)$.  By the
Riemann existence theorem this same generating triple gives a complex
$G$-cover of $\PP^1$ with the same signature, and hence, by
Riemann--Hurwitz, with the same genus.  This is the only specialization
consequence used below.

The tame Riemann--Hurwitz formula reads
\begin{equation}\label{eq:triangleRH}
  2g-2
   =|G|\left(1-\frac1a-\frac1b-\frac1c\right).
\end{equation}

\begin{remark}\label{rem:threebranch}
We shall repeatedly use the following elementary consequence of the tame
Riemann--Hurwitz formula.  Let $G\le\Aut(X)$ satisfy $p\nmid|G|$. If
\[
        \frac{|G|}{g-1}>12,
\]
then $X/G$ is rational and the cover $X\to X/G$ has exactly three branch
points.  Indeed, if $X/G$ has positive genus, Riemann--Hurwitz gives
$|G|/(g-1)\le4$.  If $X/G\simeq\PP^1$ has at least four branch points,
then $|G|/(g-1)\le12$; for four branch points the smallest positive
orbifold defect is $1/6$, attained by $(2,2,2,3)$, and for at least five
branch points the defect is at least $1/2$.
\end{remark}

All characteristic-zero actions used below have genus at most $48$.
We use Breuer's classification of finite groups of holomorphic
(equivalently, orientation-preserving) automorphisms of compact Riemann
surfaces \cite[Chapter~5 and the summary table on p.~91]{Breuer}.
We shall also use Burnside's normal $p$-complement theorem and his
$p^a q^b$ theorem in the forms stated in
\cite[Theorems~5.13 and~7.8]{Isaacs}.

\subsection{\texorpdfstring{The curves $\cH_m$}{The curves Hm}}

For later use in the prime-field case, let $m$ be a divisor of $p+1$ and put
\[
            \cH_m:\qquad y^m=x^p+x .
\]
Then
\begin{equation}\label{eq:Hmgenus}
            g(\cH_m)=\frac{(m-1)(p-1)}2.
\end{equation}
For $2\le m<p+1$ one has
\begin{equation}\label{eq:Hmaut}
            |\Aut(\cH_m)|=m\,p(p^2-1);
\end{equation}
indeed $\Aut(\cH_m)$ contains a normal cyclic group of order $m$ with
quotient $PGL(2,p)$; see \cite[Lemma~2.3]{BMT}.


\section{The tame case}\label{sec:tame}


\begin{thebibliography}{99}
\bibitem[BMT]{BMT} D.~Bartoli, M.~Montanucci and F.~Torres, \emph{$\F_{p^2}$-maximal curves with many automorphisms are Galois-covered by the Hermitian curve}, Adv. Geom. \textbf{21} (2021), 325--336.
\bibitem[Breuer]{Breuer} T.~Breuer, \emph{Characters and Automorphism Groups of Compact Riemann Surfaces}, London Math. Soc. Lecture Note Ser. \textbf{280}, Cambridge University Press, Cambridge, 2000.
\bibitem[FGT]{FGT} R.~Fuhrmann, A.~Garcia and F.~Torres, \emph{On maximal curves}, J. Number Theory \textbf{67} (1997), no.~1, 29--51.
\bibitem[GMP]{GMP} R.~Guralnick, B.~Malmskog and R.~Pries, \emph{The automorphism groups of a family of maximal curves}, J. Algebra \textbf{361} (2012), 92--106.
\bibitem[Isaacs]{Isaacs} I.~M.~Isaacs, \emph{Finite Group Theory}, Graduate Studies in Mathematics \textbf{92}, American Mathematical Society, Providence, RI, 2008.
\bibitem[KT]{KT} G.~Korchm\'aros and F.~Torres, \emph{On the genus of a maximal curve}, Math. Ann. \textbf{323} (2002), 589--608.
\bibitem[Lachaud]{Lachaud} G.~Lachaud, \emph{Sommes d'Eisenstein et nombre de points de certaines courbes alg\'ebriques sur les corps finis}, C. R. Acad. Sci. Paris S\'er. I Math. \textbf{305} (1987), no.~16, 729--732.
\bibitem[SGA1]{SGA1} A.~Grothendieck, \emph{Rev\^etements \'etales et groupe fondamental (SGA~1)}, Lecture Notes in Mathematics \textbf{224}, Springer, Berlin, 1971.
\end{thebibliography}
\end{document}
